\begin{helper}
Fix a history cell, the terminal coordinate set \(U\), the nodup terminal
reveal order, a fixed blue support label \(B\), and a fixed red label \(J\).
Use the branch/transcript interface supplied by
\(\noderef{FixedSetHistoryCellRedBranchTranscriptCells}\): branch labels
\(\beta\), transcript labels \(\tau=(P_\tau,A_\tau,Z_\tau)\), decoded red
supports \(R_{\beta,J}\), the realized-cell predicate, and the complete
assignment compatibility predicate.  Assume the fixed-\((B,J)\) realized
geometry, compatibility soundness, compatibility uniqueness, and the
non-realized cell exclusion exported by that interface.  In addition assume
the fixed residual scan/completeness interface for this \(B,J\): every
residual assignment \(A_{\rm res}\subseteq U\) satisfying
\[
P_\tau\setminus(B\cup R_{\beta,J})\subseteq A_{\rm res},\qquad
A_{\rm res}\cap(A_\tau\setminus Z_\tau)=\emptyset
\]
reconstructs to a complete assignment compatible with \((\beta,\tau)\), and
if the same residual assignment has compatible reconstructed complete
assignments for two realized pairs \((\beta,\tau)\) and
\((\beta',\tau')\), then \(\beta=\beta'\) and \(\tau=\tau'\).

Then there are fixed-\((B,J)\) released residual blocks
\[
\mathcal R_{\beta,\tau}\subseteq 2^U
\]
for every realized pair \((\beta,\tau)\), with the following properties.
Every block assignment \(A_{\rm res}\in\mathcal R_{\beta,\tau}\) is a
residual assignment on \(U\), and the reconstructed full terminal assignment
\[
(A_{\rm res}\setminus Z_\tau)\cup B\cup R_{\beta,J}
\]
is compatible with \((\beta,\tau)\).  Conversely, if a residual assignment
\(A_{\rm res}\subseteq U\) reconstructs to a full assignment compatible with
\((\beta,\tau)\), then \(A_{\rm res}\in\mathcal R_{\beta,\tau}\).  Blocks
for distinct realized pairs are disjoint.  Finally, the full residual
cylinder obtained after releasing the fixed supports and the selected
absences is contained in the block:
\[
\{A_{\rm res}\subseteq U:
P_\tau\setminus(B\cup R_{\beta,J})\subseteq A_{\rm res},\
A_{\rm res}\cap(A_\tau\setminus Z_\tau)=\emptyset\}
\subseteq \mathcal R_{\beta,\tau}.
\]
\end{helper}
\begin{proof}
The construction is the fixed-\((B,J)\) residual version of the reveal-order
argument in the paper.  The paper reveals all blue terminal choices first
and then the red choices; after the blue trace is fixed, the red choices
available for \(U_R\) are determined by that trace.  The node
\(\noderef{FixedSetHistoryCellRedBranchTranscriptCells}\) is precisely the
finite branch/transcript interface for this deterministic scan: \(\beta\)
records the complete blue terminal trace, \(J\) records the fixed red support
choice, and \(\tau=(P_\tau,A_\tau,Z_\tau)\) records the prefix transcript
and selected absences.

For a realized pair \((\beta,\tau)\), define the released block
\(\mathcal R_{\beta,\tau}\) to be the set of residual assignments
\(A_{\rm res}\subseteq U\) whose reconstructed full terminal assignment
\[
A^{\mathrm{full}}_{\beta,\tau}(A_{\rm res})
  =(A_{\rm res}\setminus Z_\tau)\cup B\cup R_{\beta,J}
\]
is accepted by the same fixed scan as compatible with
\((\beta,\tau)\).  This definition immediately gives the block-membership
soundness assertion: membership says \(A_{\rm res}\subseteq U\) and that the
displayed reconstruction is compatible.  The converse membership assertion
is the reverse implication in the same definition.

It remains to check cylinder coverage and disjointness.  Let
\(A_{\rm res}\) lie in the displayed residual cylinder.  The realized
geometry from
\(\noderef{FixedSetHistoryCellRedBranchTranscriptCells}\) gives
\[
B\cup R_{\beta,J}\subseteq P_\tau,\qquad
Z_\tau\subseteq A_\tau,\qquad P_\tau\cap A_\tau=\emptyset .
\]
Thus the reconstruction contains every coordinate in \(P_\tau\): the fixed
support coordinates \(B\cup R_{\beta,J}\) are inserted explicitly, and the
remaining coordinates \(P_\tau\setminus(B\cup R_{\beta,J})\) are present in
\(A_{\rm res}\).  It also avoids every coordinate of \(A_\tau\setminus
Z_\tau\), because those coordinates are excluded from \(A_{\rm res}\) by
the residual cylinder and are not reinserted.  The residual
scan/completeness hypothesis is exactly the converse implication for this
calculation: with \(B,J,\beta,\tau\) fixed, these residual cylinder
constraints imply that the reconstructed complete assignment is compatible
with \((\beta,\tau)\).  This is the formal version of the paper's
blue-first/red-second reveal order: once the complete blue trace is fixed,
the later red scan sees a determined set of relevant terminal pairs.  The
scan stability calculation is the content isolated in
\(\noderef{FixedSetHistoryCellRISIResidualScanStability}\).  Hence
\(A_{\rm res}\in\mathcal R_{\beta,\tau}\).

For disjointness, suppose one residual assignment belonged to both
\(\mathcal R_{\beta,\tau}\) and \(\mathcal R_{\beta',\tau'}\).  The two
memberships say that the appropriate reconstructed full assignments are
accepted by the fixed deterministic residual scan for the two realized
pairs.  The single-owner hypothesis is formulated at the residual assignment
level, so it applies even though the two reconstructed complete assignments
may use different selected sets or red supports.  It forces
\(\beta=\beta'\) and \(\tau=\tau'\).  This is the residual analogue of the
functional-scan uniqueness recorded in
\(\noderef{FixedSetHistoryCellRedFunctionalScanDisjointness}\), but with the
extra converse reconstruction information supplied explicitly.  Hence
blocks indexed by distinct realized pairs are disjoint.

The construction uses exactly the support demanded by the product-tree
packing step: block membership, reconstructed-assignment soundness,
converse membership from reconstructed compatibility, disjointness, and
residual-cylinder containment.  These are the fixed-\((B,J)\) residual
blocks needed for the partition and decision-tree package.
\end{proof}
